% !TEX root =  mosaic.tex
\section{Conclusion}
\label{sec:conclusion}

We presented \name, a unified analytical framework for GPU execution-space optimization. \name selects tile sizes and thread-block dimensions for a transformed program whose schedule is already fixed, reasoning over abstract shared-memory placement scenarios. Its two-stage analytical selection preserves the quantized primary optimum and uses a relaxed-cost tie-break only among tied legal discrete configurations, including across modeled placement scenarios, so \name selects these parameters without generating and evaluating large numbers of GPU kernels. Experimental results include a cross-generation PolyBench evaluation on three NVIDIA GPUs, together with deep learning operators and end-to-end workloads on the RTX 4090. Across these experiments, \name achieves geometric-mean speedups over PPCG, lower aggregate latency relative to EATSS, and slightly higher end-to-end performance than TVM MetaSchedule at substantially lower optimization overhead, suggesting that architecture-aware execution-space optimization can be addressed analytically rather than through measurement-driven search.

\name operates after transformation scheduling: it preserves the upstream schedule and mapping structure, selects tile and thread-block parameters, and leaves concrete shared-memory staging to the backend. Its cost model uses the loop and access metadata of \autoref{sec:metadata} and covers conventional CUDA-core SIMT execution, as evaluated through the PPCG and non-tensorized TVM paths. It stops there: tensor-core operand staging and MMA shapes, asynchronous copy and the pipelines built on it, shared-memory swizzling, warp specialization, and thread-block clusters lie outside it, and where those mechanisms decide performance, the parameters \name selects are not the ones that decide it. What would carry over to such a target is the structure rather than the contents of the formulation: wave quantization follows from an SM hosting an integer number of blocks, and residency from the register and shared-memory budgets, neither of which depends on what issues the arithmetic. The terms would change: \(\phi\) would normalize operand reads by MMA tiles rather than by per-thread FMAs, \(\hat R\) would count accumulator fragments, and the \(\Lambda\) family would need constants for the asynchronous paths. Irregular accesses need metadata this interface does not carry.
